\subsection{A fresh four-object census at sixteen arrows}

\begin{proposition}[Proved by complete finite computation]
\label{CP2:A-baseline-four-census}
No category on four objects with at most sixteen arrows has
classifying space of type $\mathbb{CP}^2$.
\end{proposition}
\begin{proof}
The excess-eight theorem excludes at most fifteen arrows.
At sixteen arrows the ninth-excess support theorem requires
acyclic support, and Proposition~\ref{CP2:A-baseline-four-large}
excludes every corner of order at least eight.  It remains
to check exactly sixteen arrows, four objects, acyclic
support, and corner orders at most seven.

\paragraph{Complete profile scope.}
The new program \path{converse/scope17/four16_scope.py}
enumerates all subsets of the six upper-triangular pairs on
four objects, retaining exactly the transitive connected
supports.  There are eighteen with these topological labels.
Every finite poset admits a linear extension, so no support
is lost.  It assigns positive integers to the four diagonal
positions and every supported cross-hom, with total sixteen,
and restricts the diagonal entries to at most seven.  This
gives $103370$ labelled decorations.  For each it tries every
object permutation and the opposite matrix, retaining the
least upper-triangular representative.  The resulting
$18217$ profiles cover every category in the stated scope;
the use of opposition preserves the homotopy type of the
nerve.

The following sufficient obstructions remove $17935$ of
these profiles.  Several can apply to the same profile;
the ledger records all reasons, and does not add overlapping
counts.  Singleton cross-homs use
Lemma~\ref{CP2:A-singleton-support}; four-object support has
no fourth homology.  Semisimple corners of orders one or
two use the support-height bound.  A full singleton cut is
a join.  A cut into two blocks with zero second cohomology
uses Lemma~\ref{CP2:A-two-block}; each proper block has at
most three objects.  The block test proves zero second
cohomology component by component using a corner of order
at most four, the H2-cost bound, singleton cross-homs, or
semisimple corners and support height at most one.

A singleton arrow into an extreme corner invokes the
fixed-arrow corollary whenever the complementary block
passes that test.  Extreme corners of order at most four
whose every stalk has size one or two use the binary-stalk
lemma, under the necessary rank-one cohomology hypothesis.
With at most two parallel excesses and corners at most
four, Lemma~\ref{CP2:A-high-two-parallel} applies.  More
generally, if no two parallel hom-sets are serial and every
corner of potentially nonzero second cohomology is incident
to a singleton edge, every global second class vanishes on
those corners: the singleton-edge join makes its endpoint
inclusions null-homotopic.  Thus
Lemma~\ref{CP2:A-high-parallel-cut} applies.  All support
cycle ranks here are at most three, within its hypothesis.

The scope also uses a literal singleton-comma deletion.
If an object has exactly one incoming cross arrow from a
trivial corner, its comma category from the deletion is a
point; the outgoing assertion is dual.  This reduction is
used only when the deletion saves at least three arrows.
The remaining three-object category then has at most
thirteen arrows and is excluded by the unrestricted
excess-eight theorem.  No unproved sixteen-arrow baseline
is invoked by this cut.

For a support tree, transitivity implies that there is no
chain of two strict comparisons.  Consequently its nerve
is the union of its edge-collage nerves along their corner
nerves.  The tree test is used only when every internal
corner has order at most five.  Such corner spaces have
zero rational homology in degrees at least three.  The
iterated Mayer--Vietoris sequences therefore pass bounded
homology to every edge and make fourth-cohomology
restriction to the edges injective.  Every edge is
required to satisfy one of the following already proved
tests: at most two cross arrows; both corners of order at
most four; at most eleven total arrows; one corner of order
at most two and the other of order at most eight; or corner
orders three and five.  The fourth case uses exact orbit
comparison and the bounded order-eight act theorem; the
last uses the $P_5$ graph theorem, after the small-corner
classification.  Thus every restriction of a global
degree-two square is zero, and injectivity makes the
global square zero.

Finally, a star with center order at most eight and all
leaves of order at most two admits exact orbit comparison
at every leaf.  This is a rational homology and cohomology
ring comparison to a star with trivial leaves.  Identifying
those leaf points adds only circles to its classifying
space and gives an ordinary act cone at the center.
Boundedness and all cohomology in degrees at least two
are preserved, so the order-eight act theorem applies.
These are all the profile cuts in the code.  An independent
targeted review checked every predicate against the named
lemma, including the transfer and deletion hypotheses.

\paragraph{Complete table search in the residual scope.}
There are $282$ residual profiles.  The largest normalized
two-chain basis has $66$ elements.  The program
\path{converse/scope17/four16_search.cpp} uses all ordinary
corner monoids from Theorem~\ref{CP2:A-monoid-census}, with
both isomorphism orientations retained.  It fixes these
corner multiplications and the identity laws.  Each other
composable pair has as its domain the complete cross-hom
of the required type.  Propagation uses only equations
$(ab)c=a(bc)$ with known inner products; it either rejects
unequal known values or forces an unknown outer product.
It then branches over every possible value of a remaining
entry.  Induction on this finite search tree proves that
every associative table with the chosen corners occurs.

The sole table-level structural cut is when all four
corner algebras are semisimple.  Nondegeneracy modulo
$32003$ of the regular trace pairing certifies
nondegeneracy over $\mathbb Q$ and hence semisimplicity.
The height-three bar bound then gives $H_4=0$.
No inherited act-cone or projective-endpoint table pruning
is enabled.  Retraction pruning in the generic engine
cannot fire, since the profiles have no arrows in both
directions between distinct objects.

The remaining tables undergo the necessary tests
$H_1(-;\mathbb F_2)=0$ and $\dim H_2(-;\mathbb F_2)=1$.
For tables passing both, exact cross-hom relabelling keys
remove only repeated typed composition tables.  Every
nonidentity composable product is in the key, while
identity products are already fixed.  Cache eviction can
retain duplicates and does not discard a new orbit.
The reported retained count is therefore a record count,
not an isomorphism-class count.

The producer supplies the full typed table, a nonboundary
two-cocycle, and either a literal three-cochain whose
coboundary is its Alexander--Whitney square, or a claim of
positive third mod-two homology.  The separate checker
\path{converse/scope17/check_four16.cpp} imports neither
the search nor the cup producer.  It reconstructs all
normalized chains, checks types, identities, every
associativity equation, connectedness, both preliminary
homology tests, the cocycle and its nonboundary class,
and every square-primitive equation.  For the other
reason it recomputes $H_3$.  The gzip end and CRC and the
generator/checker record totals are also checked.

All $282$ profile intervals completed.  The producer visited
$155765567$ corner-representative associative tables, of which
$15341745$ passed the first two homology tests before relabelling.
The independent checker accepted $1721983$ retained records:
$1719138$ have literal zero-square primitives and $2845$ have
positive third mod-two homology.  No record remains unexcluded.
The complete join \path{outputs/four16/scope_join.json} verifies
all profile indices, the exact residual ledger, every completed
interval, generator/checker count equality, source hashes and
all closed certificate-stream hashes.  The proof ledger and all
$282$ checked streams are retained under \path{outputs/four16/}.
The earlier interrupted FIFO-transport attempt is preserved
separately and contributes no completion claim.

This completes the residual exclusion by
Proposition~\ref{CP2:A-cup-certificate} and hence proves
the proposition in its full stated scope.
\end{proof}
